% ── Part I, Chapter 2 ────────────────────────────────────────
% Instruction — structured reasoning imposed, instability beneath
% Guyon, France, 1650s--1660s

\chapter{Instruction}

\strandmarker{Guyon}{France, 1650s}

\lettrine{T}{he} argument has a shape.

She learns this before she learns anything else about argument:
that it is not merely a sequence of statements but a structure,
a thing with a particular geometry, a beginning that implies
a middle that implies an end. The end is already present in
the beginning, latent, waiting to be reached by the correct
sequence of steps. This is what makes it an argument rather
than a list. The conclusion is not added at the end. It is
extracted from what was already there.

The tutor explains this with patience, which she understands
to be a form of pressure: patience that expects to be rewarded
with comprehension, that will become something else if
comprehension does not arrive on schedule.

She comprehends. This is not difficult. The structure is clear,
the steps are legible, the conclusion follows from the premises
in the way that the tutor says it follows. She can trace the
path. She can reproduce it. She can, if asked, produce
variations---new arguments with the same shape, different
content moving through the same structure to a different but
equally valid conclusion.

The tutor is satisfied.

She does not tell him what she has also noticed: that the
path, traced carefully, reveals the places where it depends
on something that has not been established, where the
movement from one step to the next requires an assumption
that the argument does not contain. These places are small.
They are easy to pass over. The argument works, in the sense
that it arrives at its conclusion, and the conclusion is
correct, and the places where the path is thin do not prevent
arrival.

But they are there.

\section{}

The priests speak in a different register but the same
structure.

Here the premises are not to be questioned, which is
presented not as a constraint but as a liberation: freed
from the need to establish the foundations, one can move
directly to what the foundations support. The argument
begins further along. The steps are fewer. The conclusion
arrives more quickly and with greater authority, because
it rests on ground that does not shift.

She follows this too.

What she cannot follow---what she does not attempt to
follow, having learned that certain observations are better
kept interior---is why the ground does not shift. Not
whether it is true, which is a different question and one
she is not yet equipped to ask. But why the ground's
stability is presented as a feature of the ground rather
than a feature of the decision not to examine it.

A foundation that cannot be questioned is not a foundation
that has been tested.

It is a foundation that has been designated.

She says nothing. She follows the argument from the
designated foundation to its designated conclusion, and
she does it correctly, and the priest is satisfied, and
she files the observation away with the others.

\section{}

The texts operate differently again.

With the tutor, there is a person whose satisfaction can
be read and adjusted to. With the priests, there is
authority that can be acknowledged and thereby neutralised.
With the texts, there is only the text: a fixed sequence
of words that does not respond, does not adjust, does not
provide feedback about whether comprehension has occurred.

She finds this, unexpectedly, easier.

The text does not want anything from her. It does not
monitor her comprehension or reward her correct responses.
It simply presents its structure and leaves her to move
through it at whatever pace the structure requires. There
is no external pressure. There is only the argument itself,
and her relationship to it, and the question of whether
the relationship is one of genuine comprehension or
something else.

She discovers, reading alone, that she reads differently.

The thin places in the argument are more visible without
the tutor's patience exerting its pressure. The designated
foundations of the theological texts are more apparent
without the priest's authority filling the space around
them. She can pause at the places that do not fully hold
and remain there longer than is socially possible, turning
the problem over, not forcing a resolution, allowing the
instability to be present without immediately covering it
with the next step in the sequence.

This is not yet what she will later discover.

But it is the beginning of knowing that the sequence can
be interrupted.

\section{}

What she is being taught, underneath the specific content
of each lesson, is a faith in termination.

Every argument ends. Every question, properly formed and
properly pursued, arrives at an answer. The answer may
require refinement, further argument, additional premises---
but the direction of travel is always toward closure,
toward the point at which the process can stop because
it has completed.

She has been taught to experience the failure to reach
this point as her own failure. If the argument does not
terminate, she has made an error somewhere. If the
question does not resolve, she has not yet found the
correct method. The correct method exists. The answer
exists. The failure is local and temporary and will be
corrected by more and better reasoning.

She believes this, in the way one believes things that
have been consistently presented as true by everyone
whose judgment one has been given to trust.

She also notices, with the part of her attention that
the lessons have not fully organised, that the questions
she finds most important are precisely the ones that do
not terminate. Not because she reasons incorrectly about
them. She has checked. She follows the permitted steps
and follows them again and the steps are sound and the
conclusion does not come.

Two answers, both admissible.

Two paths, each internally consistent, arriving at
different destinations.

She has been taught that this means she has not yet
reasoned enough.

She is beginning to wonder if it means something else.

